\begin{table}[H]
\centering
\caption{Description of variables -- Controls used}
\begin{tabular}{|p{0.35\textwidth}|p{0.55\textwidth}|}
\hline
\textbf{Variables} & \textbf{Definition} \\
\hline

\texttt{Number of rooms} & Self-explanatory \\
\hline

\texttt{Average room size} & Self-explanatory \\
\hline

\texttt{Floor} & Number of floors \\
\hline

\texttt{Age} & 
\begin{tabular}[t]{l}
- Age $< 3$ \\
- After the Tohoku earthquake (from 2011) \\
- After the Kobe earthquake (1995--2010) \\
- After the para-seismic law (1981--1994) \\
- After the City Planning Law (1968--1980) \\
- After the WW2 (1945--1967)
\end{tabular}
\\
\hline

\texttt{Building structure} &
\begin{tabular}[t]{l}
- Wooden structure \\
- Block structure \& Reinforced concrete block \\
- Steel frame structure \& Lightweight steel frame \\
- Reinforced concrete \& Precast concrete \& High-performance precast concrete \& Alkaline concrete \\
- Steel reinforced concrete
\end{tabular}
\\
\hline

\texttt{Distance to station} & Distance to the nearest station in Km. In the Japanese case, stations are a good proxy for city centers. \\
\hline

\texttt{Density} & 
Density Inhabited Districts (DID) is defined by the Population Census, where areas with a population density of 4,000 people per square kilometer or higher are adjacent to each other and have a total population of 5,000 or more.

\begin{itemize}
\setlength{\itemsep}{0pt}
\item 0 if DID is not defined
\item $\log\left(\frac{\text{Number of people living in DID}}{\text{Size of the DID area}}\right)$ if DID is defined
\end{itemize}
\\
\hline

\texttt{Net migration rate} & Number of in-migrants over number of total population in 2015. \\
\hline

\end{tabular}
\end{table}